% =====================================================================
%  2bat-ud6-7.tex  —  SESSIÓ 7: El marge d'error i la mida de la mostra
%  (sense preàmbul: s'inclou des de main.tex)
% =====================================================================
\horatitol{Sessió 7 — El marge d'error i la mida de la mostra}%
{Treball tecnològic: explorar amb el full de càlcul com el marge d'error depèn de la mida de la mostra, i el «cost de la precisió».}

\subsection{Idea clau}

Amb el \textbf{full de càlcul}, explorem una relació clau que sovint es malinterpreta:
\textbf{com depèn el marge d'error de la mida de la mostra}. Comprovareu que més mostra
dona menys marge (més precisió), però que la millora \textbf{no és proporcional}: a partir
d'un cert punt, augmentar la mostra surt molt car per a una millora petita.

\begin{idea}
\textbf{Marge d'error i mida de la mostra.}
\begin{itemize}[leftmargin=1.4em, itemsep=2pt]
  \item Com més \textbf{gran} és la mostra ($n$), més \textbf{petit} és el marge d'error:
        l'estimació és més precisa.
  \item Però la relació \textbf{no és proporcional}: per \textbf{reduir el marge a la
        meitat} cal \textbf{multiplicar la mostra per 4}, no per 2.
  \item Per tant, a partir d'un cert punt, guanyar una mica de precisió costa moltíssimes
        enquestes més: és el \textbf{cost de la precisió}.
\end{itemize}
\textbf{Al full de càlcul} es pot calcular el marge aproximat (per a un resultat del
$50\%$, confiança $95\%$) amb una fórmula del tipus $\texttt{=0,98/ARREL(n)}$ i veure com
canvia en variar $n$.
\end{idea}

\subsection{Exemples}

\begin{exemple}[la taula del marge segons la mida]
Amb el full de càlcul obtenim el marge d'error aproximat per a diferents mides de mostra:
\begin{center}
\renewcommand{\arraystretch}{1.2}
\begin{tabular}{c c}
\toprule
\textbf{Mida de la mostra $n$} & \textbf{Marge d'error aprox.} \\
\midrule
100  & $\pm 9,8\%$ \\
400  & $\pm 4,9\%$ \\
900  & $\pm 3,3\%$ \\
1\,600 & $\pm 2,5\%$ \\
2\,500 & $\pm 2,0\%$ \\
\bottomrule
\end{tabular}
\end{center}
Fixa't: de $100$ a $400$ (×4) el marge passa de $\approx 9,8\%$ a $\approx 4,9\%$ (la
meitat). Cal \textbf{quadruplicar} la mostra per \textbf{reduir el marge a la meitat}.
\end{exemple}

\begin{exemple}[el cost de la precisió, en un gràfic]
\begin{center}
\begin{tikzpicture}
\begin{axis}[udaxis, width=10.5cm, height=5cm,
    xlabel={\small mida de la mostra $n$}, ylabel={\small marge (\%)},
    xmin=0, xmax=2600, ymin=0, ymax=11, xtick={100,400,900,1600,2500}, ytick={2,4,6,8,10},
    x tick label style={font=\scriptsize, rotate=40, anchor=east}]
  \addplot[udblue, thick, mark=*, mark size=1.4pt] coordinates
    {(100,9.8)(400,4.9)(900,3.27)(1600,2.45)(2500,1.96)};
\end{axis}
\end{tikzpicture}

\figcap{El marge cau de pressa al principi i després molt a poc a poc: a partir d'un cert punt, més mostra amb prou feines millora la precisió.}
\end{center}
La corba baixa molt entre $100$ i $400$, però a partir de $900$ gairebé es fa plana:
duplicar la mostra ja no compensa.
\end{exemple}

\subsection{Exercicis resolts}

\begin{exresolt}[interpretar la taula]
Segons la taula, per passar d'un marge de $\pm 4,9\%$ a un de $\pm 2,5\%$ (aproximadament la
meitat), com ha de canviar la mostra?
\solucio
De $\pm 4,9\%$ (a $n=400$) a $\pm 2,5\%$ (a $n=1\,600$): la mostra passa de $400$ a
$1\,600$, és a dir, es \textbf{multiplica per 4}. Per reduir el marge a la meitat cal
quadruplicar la mida.
\resposta{Cal multiplicar la mostra per $4$ (de $400$ a $1\,600$).}
\end{exresolt}

\begin{exresolt}[decisió pràctica]
Una entitat té una enquesta amb $n=900$ (marge $\pm 3,3\%$) i es planteja passar a
$n=2\,500$ (marge $\pm 2,0\%$). Val la pena? Raona-ho.
\solucio
Passar de $900$ a $2\,500$ gairebé \textbf{triplica} el nombre d'enquestes (i el cost), i el
marge només baixa d'un $\pm 3,3\%$ a un $\pm 2,0\%$ (una millora d'$1,3$ punts). Segons el
pressupost, probablement \textbf{no compensa}: per a la majoria de decisions socials, un
marge de $\pm 3,3\%$ ja és prou bo.
\resposta{Sovint no compensa: molt més cost per una millora petita del marge.}
\end{exresolt}

\subsection{Exercicis per fer}

\noindent\textit{Fes servir la taula de l'exemple (marge segons $n$).}

\begin{exercicis}
\begin{exlist}
  \item Quin marge d'error aproximat té una enquesta amb $n=400$? I amb $n=1\,600$?
  \item Per reduir el marge de $\pm 9,8\%$ (a $n=100$) a la meitat, quina mida de mostra
        cal? Quin factor és respecte de $100$?
  \item Ordena de \emph{més} a \emph{menys} precisa aquestes enquestes segons la mida:
        $n=100$, $n=900$, $n=2\,500$.
  \item Explica amb les teves paraules per què «doblar la mostra no redueix el marge a la
        meitat».
  \item \textbf{(Decisió)} Una enquesta amb $n=1\,600$ té un marge de $\pm 2,5\%$. Passar a
        $n=2\,500$ (marge $\pm 2,0\%$) suposa moltes més enquestes. Val la pena per guanyar
        mig punt de precisió? Raona segons el context.
  \item \textbf{(Interpretació)} Per què les enquestes reals solen fer servir mostres d'uns
        pocs milers de persones i no de centenars de milers, encara que la població sigui
        enorme?
\end{exlist}
\end{exercicis}

% ---------------------------------------------------------------------
\fullsolucions{Sessió 7}
\begin{solucions}
\begin{sollist}
  \item $n=400$: $\pm 4,9\%$; $n=1\,600$: $\pm 2,5\%$.
  \item La meitat de $\pm 9,8\%$ és $\pm 4,9\%$, que correspon a $n=400$: cal
        \textbf{multiplicar la mostra per 4} (de $100$ a $400$).
  \item De més a menys precisa: $n=2\,500$ ($\pm 2,0\%$), $n=900$ ($\pm 3,3\%$), $n=100$
        ($\pm 9,8\%$).
  \item Resposta oberta. Idea: el marge depèn de l'arrel de la mostra, no de la mostra
        directament; per això cal quadruplicar $n$ per reduir el marge a la meitat, i
        doblar-la només el redueix una mica.
  \item Resposta oberta. Idea: sovint no compensa: gairebé duplicar les enquestes (i el
        cost) per guanyar mig punt de precisió rarament es justifica, tret que la decisió
        exigeixi molta exactitud.
  \item Resposta oberta. Idea: perquè, a partir d'uns quants milers, el marge ja és petit i
        augmentar més la mostra amb prou feines el millora, mentre que el cost creix molt;
        s'escull la mida que equilibra precisió i cost.
\end{sollist}
\end{solucions}
